% --- Archon physics formalization source begin ---
% archon:physics
% archon:covers PhyXMiniProblems/problem_phyx_mini_0370.lean
% archon:source-report reports/phyx_mini/problem_phyx_mini_0370.source.json

\chapter{Physics problem phyx\_mini\_0370}
\label{ch:PhyXMiniProblems_problem_phyx_mini_0370}

\paragraph{Problem source.}
\#\# Physical scenario

The figure shows two different processes by which 80 mol of gas move from state 1 to state 2. The dashed line is an isotherm.

\#\# Question

What maximum temperature is reached along the straight-line process?

\#\# Answer choices

- A: A: 338K

- B: B: 556K

- C: C: 636K

- D: D: 784K

\#\# Figure caption (auxiliary; use the image as primary evidence)

The image is a pressure-volume (\textbackslash{}( p \textbackslash{})-\textbackslash{}( V \textbackslash{})) diagram labeled with an isotherm. 

- **Axes**: 
  - The vertical axis represents pressure (\textbackslash{}( p \textbackslash{})) measured in kilopascals (kPa) ranging from 0 to 200.
  - The horizontal axis represents volume (\textbackslash{}( V \textbackslash{})) measured in cubic meters (m\textbackslash{}(\textasciicircum{}3\textbackslash{})) ranging from 0 to 3.

- **Curve**:
  - A red dashed curve labeled "Isotherm" extends from the top left to the bottom right, showing how pressure changes with volume at a constant temperature.

- **Points**:
  - Two primary points are plotted: 
    - Point 1 at approximately \textbackslash{}( V = 2 \textbackslash{}, \textbackslash{}text\{m\}\textasciicircum{}3 \textbackslash{}), \textbackslash{}( p = 110 \textbackslash{}, \textbackslash{}text\{kPa\} \textbackslash{})
    - Point 2 at approximately \textbackslash{}( V = 1 \textbackslash{}, \textbackslash{}text\{m\}\textasciicircum{}3 \textbackslash{}), \textbackslash{}( p = 220 \textbackslash{}, \textbackslash{}text\{kPa\} \textbackslash{})

- **Arrows**:
  - Two black arrows point from point 2 to point 1 along the curve, indicating a process or transition occurring from higher to lower pressure and volume.

\#\# Dataset metadata

Ideal Gases and Kinetic Theory; Physical Model Grounding Reasoning, Numerical Reasoning

\paragraph{Recorded answer/context.}
A: A: 338K

\paragraph{Figure/image path.}
/root/proposal\_for\_physic/science-mango/phyx\_mini\_run/phyx\_data/test\_image/370.png

\paragraph{Formalization target.}
create a compiling Lean file with sorry bodies at `PhyXMiniProblems/problem\_phyx\_mini\_0370.lean`. The Lean declarations must preserve the physical quantities, dimensions or dimensional roles, figure labels, governing-law hypotheses, and final relation expressed by this problem.
Use Mathlib/Physlib names found through LeanExplore where available. If a physics API is missing, introduce faithful local abstractions rather than scalar placeholder aliases.

\begin{theorem}[Physics formalization target]
\label{thm:physics:phyx_mini_0370:target}
\lean{PhyXMiniProblems.ProblemPhyXMini0370.maximum_temperature_on_straight_line_is_answer_A}
\uses{def:physics:phyx-mini-0370:phyxminiproblems-problemphyxmini0370-amountofsubstancescale, def:physics:phyx-mini-0370:phyxminiproblems-problemphyxmini0370-idealgastwopathsetup, def:physics:phyx-mini-0370:phyxminiproblems-problemphyxmini0370-matchesproblemdescription, def:physics:phyx-mini-0370:phyxminiproblems-problemphyxmini0370-matchesprimarypressurevolumediagram, def:physics:phyx-mini-0370:phyxminiproblems-problemphyxmini0370-hasphysicalprocessparameters, def:physics:phyx-mini-0370:phyxminiproblems-problemphyxmini0370-satisfiesidealgasandprocesslaws, def:physics:phyx-mini-0370:phyxminiproblems-problemphyxmini0370-usestextbookmolargasconstant, def:physics:phyx-mini-0370:phyxminiproblems-problemphyxmini0370-straightlinetemperatureinkelvin, def:physics:phyx-mini-0370:phyxminiproblems-problemphyxmini0370-ismaximumstraightlinetemperatureat, def:physics:phyx-mini-0370:phyxminiproblems-problemphyxmini0370-recordeddatasetanswer, def:physics:phyx-mini-0370:phyxminiproblems-problemphyxmini0370-roundstodisplayedwholekelvin, def:physics:phyx-mini-0370:phyxminiproblems-problemphyxmini0370-isuniqueclosestdisplayedtemperature}
The assigned autoformalize agent should translate this physics problem into Lean declarations in the covered file, with theorem and lemma proof bodies written as `by sorry`.
\end{theorem}
\begin{proof}
This is an autoformalization task, not a proof task. Produce statements that can later be proved by the physics prover without weakening the physical model.
\end{proof}

% --- Archon declaration topology begin ---
\section{Declaration topology}
\begin{definition}[Volume Quantity]
  \label{def:physics:phyx-mini-0370:phyxminiproblems-problemphyxmini0370-volumequantity}
  \lean{PhyXMiniProblems.ProblemPhyXMini0370.VolumeQuantity}
  A nonnegative physical volume with length-cubed dimension
\end{definition}
\begin{proof}
  This is the declared model definition of Volume Quantity.
\end{proof}
\begin{definition}[Amount Of Substance Scale]
  \label{def:physics:phyx-mini-0370:phyxminiproblems-problemphyxmini0370-amountofsubstancescale}
  \lean{PhyXMiniProblems.ProblemPhyXMini0370.AmountOfSubstanceScale}
  Physlib does not currently provide an amount-of-substance base dimension. This interface keeps the gas amount abstract and exposes only its calibrated mole readout
\end{definition}
\begin{proof}
  This is the declared model definition of Amount Of Substance Scale.
\end{proof}
\begin{definition}[volume In Cubic Meters]
  \label{def:physics:phyx-mini-0370:phyxminiproblems-problemphyxmini0370-volumeincubicmeters}
  \lean{PhyXMiniProblems.ProblemPhyXMini0370.volumeInCubicMeters}
  \uses{def:physics:phyx-mini-0370:phyxminiproblems-problemphyxmini0370-volumequantity}
  Read a physical volume in coherent SI cubic metres
\end{definition}
\begin{proof}
  This is the declared model definition of volume In Cubic Meters.
\end{proof}
\begin{definition}[pressure In Pascals]
  \label{def:physics:phyx-mini-0370:phyxminiproblems-problemphyxmini0370-pressureinpascals}
  \lean{PhyXMiniProblems.ProblemPhyXMini0370.pressureInPascals}
  Read a physical pressure in coherent SI pascals
\end{definition}
\begin{proof}
  This is the declared model definition of pressure In Pascals.
\end{proof}
\begin{definition}[pressure In Kilopascals]
  \label{def:physics:phyx-mini-0370:phyxminiproblems-problemphyxmini0370-pressureinkilopascals}
  \lean{PhyXMiniProblems.ProblemPhyXMini0370.pressureInKilopascals}
  \uses{def:physics:phyx-mini-0370:phyxminiproblems-problemphyxmini0370-pressureinpascals}
  Read a physical pressure in the kilopascals printed on the diagram
\end{definition}
\begin{proof}
  This is the declared model definition of pressure In Kilopascals.
\end{proof}
\begin{definition}[temperature In Kelvin]
  \label{def:physics:phyx-mini-0370:phyxminiproblems-problemphyxmini0370-temperatureinkelvin}
  \lean{PhyXMiniProblems.ProblemPhyXMini0370.temperatureInKelvin}
  Convert a stored Physlib absolute temperature to a kelvin readout
\end{definition}
\begin{proof}
  This is the declared model definition of temperature In Kelvin.
\end{proof}
\begin{definition}[Gas Model]
  \label{def:physics:phyx-mini-0370:phyxminiproblems-problemphyxmini0370-gasmodel}
  \lean{PhyXMiniProblems.ProblemPhyXMini0370.GasModel}
  The equation-of-state model named by the problem
\end{definition}
\begin{proof}
  This is the declared model definition of Gas Model.
\end{proof}
\begin{definition}[Endpoint]
  \label{def:physics:phyx-mini-0370:phyxminiproblems-problemphyxmini0370-endpoint}
  \lean{PhyXMiniProblems.ProblemPhyXMini0370.Endpoint}
  The two black endpoint labels in the primary image
\end{definition}
\begin{proof}
  This is the declared model definition of Endpoint.
\end{proof}
\begin{definition}[Process Path]
  \label{def:physics:phyx-mini-0370:phyxminiproblems-problemphyxmini0370-processpath}
  \lean{PhyXMiniProblems.ProblemPhyXMini0370.ProcessPath}
  The two alternative paths connecting the endpoint states
\end{definition}
\begin{proof}
  This is the declared model definition of Process Path.
\end{proof}
\begin{definition}[Axis Quantity]
  \label{def:physics:phyx-mini-0370:phyxminiproblems-problemphyxmini0370-axisquantity}
  \lean{PhyXMiniProblems.ProblemPhyXMini0370.AxisQuantity}
  Physical quantity represented by a diagram axis
\end{definition}
\begin{proof}
  This is the declared model definition of Axis Quantity.
\end{proof}
\begin{definition}[Axis Unit]
  \label{def:physics:phyx-mini-0370:phyxminiproblems-problemphyxmini0370-axisunit}
  \lean{PhyXMiniProblems.ProblemPhyXMini0370.AxisUnit}
  Unit printed next to a diagram axis
\end{definition}
\begin{proof}
  This is the declared model definition of Axis Unit.
\end{proof}
\begin{definition}[Path Style]
  \label{def:physics:phyx-mini-0370:phyxminiproblems-problemphyxmini0370-pathstyle}
  \lean{PhyXMiniProblems.ProblemPhyXMini0370.PathStyle}
  Visual style used to distinguish the two plotted processes
\end{definition}
\begin{proof}
  This is the declared model definition of Path Style.
\end{proof}
\begin{definition}[Thermodynamic State]
  \label{def:physics:phyx-mini-0370:phyxminiproblems-problemphyxmini0370-thermodynamicstate}
  \lean{PhyXMiniProblems.ProblemPhyXMini0370.ThermodynamicState}
  \uses{def:physics:phyx-mini-0370:phyxminiproblems-problemphyxmini0370-volumequantity}
  A thermodynamic equilibrium state with dimensionful observables
\end{definition}
\begin{proof}
  This is the declared model definition of Thermodynamic State.
\end{proof}
\begin{definition}[Pressure Volume Diagram]
  \label{def:physics:phyx-mini-0370:phyxminiproblems-problemphyxmini0370-pressurevolumediagram}
  \lean{PhyXMiniProblems.ProblemPhyXMini0370.PressureVolumeDiagram}
  \uses{def:physics:phyx-mini-0370:phyxminiproblems-problemphyxmini0370-endpoint, def:physics:phyx-mini-0370:phyxminiproblems-problemphyxmini0370-processpath, def:physics:phyx-mini-0370:phyxminiproblems-problemphyxmini0370-axisquantity, def:physics:phyx-mini-0370:phyxminiproblems-problemphyxmini0370-axisunit, def:physics:phyx-mini-0370:phyxminiproblems-problemphyxmini0370-pathstyle}
  Scalar fields in this structure are literal coordinate readouts from the supplied raster; physical states are stored separately in the experiment and are connected to these coordinates below
\end{definition}
\begin{proof}
  This is the declared model definition of Pressure Volume Diagram.
\end{proof}
\begin{definition}[Ideal Gas Two Path Setup]
  \label{def:physics:phyx-mini-0370:phyxminiproblems-problemphyxmini0370-idealgastwopathsetup}
  \lean{PhyXMiniProblems.ProblemPhyXMini0370.IdealGasTwoPathSetup}
  \uses{def:physics:phyx-mini-0370:phyxminiproblems-problemphyxmini0370-amountofsubstancescale, def:physics:phyx-mini-0370:phyxminiproblems-problemphyxmini0370-gasmodel, def:physics:phyx-mini-0370:phyxminiproblems-problemphyxmini0370-endpoint, def:physics:phyx-mini-0370:phyxminiproblems-problemphyxmini0370-processpath, def:physics:phyx-mini-0370:phyxminiproblems-problemphyxmini0370-thermodynamicstate, def:physics:phyx-mini-0370:phyxminiproblems-problemphyxmini0370-pressurevolumediagram}
  The physical gas sample and both processes. A path parameter of 0 corresponds to state 1, and a parameter of 1 corresponds to state 2. Only parameters in the closed unit interval describe the plotted processes
\end{definition}
\begin{proof}
  This is the declared model definition of Ideal Gas Two Path Setup.
\end{proof}
\begin{definition}[gas Amount In Moles]
  \label{def:physics:phyx-mini-0370:phyxminiproblems-problemphyxmini0370-gasamountinmoles}
  \lean{PhyXMiniProblems.ProblemPhyXMini0370.gasAmountInMoles}
  \uses{def:physics:phyx-mini-0370:phyxminiproblems-problemphyxmini0370-amountofsubstancescale, def:physics:phyx-mini-0370:phyxminiproblems-problemphyxmini0370-idealgastwopathsetup}
  Mole readout of the physical gas amount
\end{definition}
\begin{proof}
  This is the declared model definition of gas Amount In Moles.
\end{proof}
\begin{definition}[process Temperature In Kelvin]
  \label{def:physics:phyx-mini-0370:phyxminiproblems-problemphyxmini0370-processtemperatureinkelvin}
  \lean{PhyXMiniProblems.ProblemPhyXMini0370.processTemperatureInKelvin}
  \uses{def:physics:phyx-mini-0370:phyxminiproblems-problemphyxmini0370-amountofsubstancescale, def:physics:phyx-mini-0370:phyxminiproblems-problemphyxmini0370-temperatureinkelvin, def:physics:phyx-mini-0370:phyxminiproblems-problemphyxmini0370-processpath, def:physics:phyx-mini-0370:phyxminiproblems-problemphyxmini0370-idealgastwopathsetup}
  Kelvin readout at a parameter value on either process
\end{definition}
\begin{proof}
  This is the declared model definition of process Temperature In Kelvin.
\end{proof}
\begin{definition}[Matches Problem Description]
  \label{def:physics:phyx-mini-0370:phyxminiproblems-problemphyxmini0370-matchesproblemdescription}
  \lean{PhyXMiniProblems.ProblemPhyXMini0370.MatchesProblemDescription}
  \uses{def:physics:phyx-mini-0370:phyxminiproblems-problemphyxmini0370-amountofsubstancescale, def:physics:phyx-mini-0370:phyxminiproblems-problemphyxmini0370-idealgastwopathsetup, def:physics:phyx-mini-0370:phyxminiproblems-problemphyxmini0370-gasamountinmoles}
  Prose data: an 80-mole ideal gas and two named alternative processes. No maximum temperature or answer choice is fixed here
\end{definition}
\begin{proof}
  This is the declared model definition of Matches Problem Description.
\end{proof}
\begin{definition}[Matches Primary Pressure Volume Diagram]
  \label{def:physics:phyx-mini-0370:phyxminiproblems-problemphyxmini0370-matchesprimarypressurevolumediagram}
  \lean{PhyXMiniProblems.ProblemPhyXMini0370.MatchesPrimaryPressureVolumeDiagram}
  \uses{def:physics:phyx-mini-0370:phyxminiproblems-problemphyxmini0370-amountofsubstancescale, def:physics:phyx-mini-0370:phyxminiproblems-problemphyxmini0370-volumeincubicmeters, def:physics:phyx-mini-0370:phyxminiproblems-problemphyxmini0370-pressureinkilopascals, def:physics:phyx-mini-0370:phyxminiproblems-problemphyxmini0370-endpoint, def:physics:phyx-mini-0370:phyxminiproblems-problemphyxmini0370-processpath, def:physics:phyx-mini-0370:phyxminiproblems-problemphyxmini0370-idealgastwopathsetup}
  Exact information transcribed from the primary raster: axis meanings and ranges, endpoint coordinates, path styles, labels, and the arrows from state 1 to state 2. The coordinate readouts are explicitly connected to the dimensionful endpoint states, and both plotted paths share those endpoints
\end{definition}
\begin{proof}
  This is the declared model definition of Matches Primary Pressure Volume Diagram.
\end{proof}
\begin{definition}[Has Physical Process Parameters]
  \label{def:physics:phyx-mini-0370:phyxminiproblems-problemphyxmini0370-hasphysicalprocessparameters}
  \lean{PhyXMiniProblems.ProblemPhyXMini0370.HasPhysicalProcessParameters}
  \uses{def:physics:phyx-mini-0370:phyxminiproblems-problemphyxmini0370-amountofsubstancescale, def:physics:phyx-mini-0370:phyxminiproblems-problemphyxmini0370-volumeincubicmeters, def:physics:phyx-mini-0370:phyxminiproblems-problemphyxmini0370-pressureinpascals, def:physics:phyx-mini-0370:phyxminiproblems-problemphyxmini0370-idealgastwopathsetup, def:physics:phyx-mini-0370:phyxminiproblems-problemphyxmini0370-gasamountinmoles, def:physics:phyx-mini-0370:phyxminiproblems-problemphyxmini0370-processtemperatureinkelvin}
  Positivity assumptions selecting physically admissible process states
\end{definition}
\begin{proof}
  This is the declared model definition of Has Physical Process Parameters.
\end{proof}
\begin{definition}[Satisfies Ideal Gas And Process Laws]
  \label{def:physics:phyx-mini-0370:phyxminiproblems-problemphyxmini0370-satisfiesidealgasandprocesslaws}
  \lean{PhyXMiniProblems.ProblemPhyXMini0370.SatisfiesIdealGasAndProcessLaws}
  \uses{def:physics:phyx-mini-0370:phyxminiproblems-problemphyxmini0370-amountofsubstancescale, def:physics:phyx-mini-0370:phyxminiproblems-problemphyxmini0370-volumeincubicmeters, def:physics:phyx-mini-0370:phyxminiproblems-problemphyxmini0370-pressureinpascals, def:physics:phyx-mini-0370:phyxminiproblems-problemphyxmini0370-pressureinkilopascals, def:physics:phyx-mini-0370:phyxminiproblems-problemphyxmini0370-idealgastwopathsetup, def:physics:phyx-mini-0370:phyxminiproblems-problemphyxmini0370-gasamountinmoles, def:physics:phyx-mini-0370:phyxminiproblems-problemphyxmini0370-processtemperatureinkelvin}
  Macroscopic laws used by the calculation: the SI molar ideal-gas equation pV = nRT at every state of both paths; affine pressure and volume coordinates along the solid straight chord; constant absolute temperature along the dashed isotherm. These are governing laws for the two processes. In particular, none states where the straight-path maximum occurs or what its value is
\end{definition}
\begin{proof}
  This is the declared model definition of Satisfies Ideal Gas And Process Laws.
\end{proof}
\begin{definition}[Uses Textbook Molar Gas Constant]
  \label{def:physics:phyx-mini-0370:phyxminiproblems-problemphyxmini0370-usestextbookmolargasconstant}
  \lean{PhyXMiniProblems.ProblemPhyXMini0370.UsesTextbookMolarGasConstant}
  \uses{def:physics:phyx-mini-0370:phyxminiproblems-problemphyxmini0370-amountofsubstancescale, def:physics:phyx-mini-0370:phyxminiproblems-problemphyxmini0370-idealgastwopathsetup}
  The standard textbook SI calibration R = 8.314 J mol⁻¹ K⁻¹. It is kept separate from the general equation of state and contains no temperature answer
\end{definition}
\begin{proof}
  This is the declared model definition of Uses Textbook Molar Gas Constant.
\end{proof}
\begin{definition}[straight Line Temperature In Kelvin]
  \label{def:physics:phyx-mini-0370:phyxminiproblems-problemphyxmini0370-straightlinetemperatureinkelvin}
  \lean{PhyXMiniProblems.ProblemPhyXMini0370.straightLineTemperatureInKelvin}
  \uses{def:physics:phyx-mini-0370:phyxminiproblems-problemphyxmini0370-amountofsubstancescale, def:physics:phyx-mini-0370:phyxminiproblems-problemphyxmini0370-idealgastwopathsetup, def:physics:phyx-mini-0370:phyxminiproblems-problemphyxmini0370-processtemperatureinkelvin}
  Temperature readout along the solid straight-line process
\end{definition}
\begin{proof}
  This is the declared model definition of straight Line Temperature In Kelvin.
\end{proof}
\begin{definition}[Is Maximum Straight Line Temperature At]
  \label{def:physics:phyx-mini-0370:phyxminiproblems-problemphyxmini0370-ismaximumstraightlinetemperatureat}
  \lean{PhyXMiniProblems.ProblemPhyXMini0370.IsMaximumStraightLineTemperatureAt}
  \uses{def:physics:phyx-mini-0370:phyxminiproblems-problemphyxmini0370-amountofsubstancescale, def:physics:phyx-mini-0370:phyxminiproblems-problemphyxmini0370-idealgastwopathsetup, def:physics:phyx-mini-0370:phyxminiproblems-problemphyxmini0370-straightlinetemperatureinkelvin}
  A parameter attains the maximum temperature on the plotted chord
\end{definition}
\begin{proof}
  This is the declared model definition of Is Maximum Straight Line Temperature At.
\end{proof}
\begin{definition}[Answer Choice]
  \label{def:physics:phyx-mini-0370:phyxminiproblems-problemphyxmini0370-answerchoice}
  \lean{PhyXMiniProblems.ProblemPhyXMini0370.AnswerChoice}
  Labels of the four printed multiple-choice answers
\end{definition}
\begin{proof}
  This is the declared model definition of Answer Choice.
\end{proof}
\begin{definition}[displayed Temperature Kelvin]
  \label{def:physics:phyx-mini-0370:phyxminiproblems-problemphyxmini0370-displayedtemperaturekelvin}
  \lean{PhyXMiniProblems.ProblemPhyXMini0370.displayedTemperatureKelvin}
  \uses{def:physics:phyx-mini-0370:phyxminiproblems-problemphyxmini0370-answerchoice}
  Whole-kelvin temperature printed beside each answer label
\end{definition}
\begin{proof}
  This is the declared model definition of displayed Temperature Kelvin.
\end{proof}
\begin{definition}[recorded Dataset Answer]
  \label{def:physics:phyx-mini-0370:phyxminiproblems-problemphyxmini0370-recordeddatasetanswer}
  \lean{PhyXMiniProblems.ProblemPhyXMini0370.recordedDatasetAnswer}
  \uses{def:physics:phyx-mini-0370:phyxminiproblems-problemphyxmini0370-answerchoice}
  Answer label recorded in the supplied dataset metadata
\end{definition}
\begin{proof}
  This is the declared model definition of recorded Dataset Answer.
\end{proof}
\begin{definition}[Rounds To Displayed Whole Kelvin]
  \label{def:physics:phyx-mini-0370:phyxminiproblems-problemphyxmini0370-roundstodisplayedwholekelvin}
  \lean{PhyXMiniProblems.ProblemPhyXMini0370.RoundsToDisplayedWholeKelvin}
  \uses{def:physics:phyx-mini-0370:phyxminiproblems-problemphyxmini0370-answerchoice, def:physics:phyx-mini-0370:phyxminiproblems-problemphyxmini0370-displayedtemperaturekelvin}
  The physical maximum is compatible with a displayed whole kelvin
\end{definition}
\begin{proof}
  This is the declared model definition of Rounds To Displayed Whole Kelvin.
\end{proof}
\begin{definition}[Is Unique Closest Displayed Temperature]
  \label{def:physics:phyx-mini-0370:phyxminiproblems-problemphyxmini0370-isuniqueclosestdisplayedtemperature}
  \lean{PhyXMiniProblems.ProblemPhyXMini0370.IsUniqueClosestDisplayedTemperature}
  \uses{def:physics:phyx-mini-0370:phyxminiproblems-problemphyxmini0370-answerchoice, def:physics:phyx-mini-0370:phyxminiproblems-problemphyxmini0370-displayedtemperaturekelvin}
  The selected displayed temperature is closer than every alternative
\end{definition}
\begin{proof}
  This is the declared model definition of Is Unique Closest Displayed Temperature.
\end{proof}
% --- Archon declaration topology end ---
% --- Archon physics formalization source end ---
